\begin{abcliste}
\abc
\result{\text{Schritt 1 ist falsch}} (und mit ihm die Schritte 4 und 5): Es gibt keine „Bewegungskraft“. Jede Kraft kommt von einem Körper, der drückt oder zieht (oder von der Erde). Sobald die Kiste Jans Hand verlassen hat, drückt nichts mehr auf sie in Bewegungsrichtung. Der Stoss ist kein Vorrat an Kraft, den die Kiste mitnimmt: Sie gleitet aus Trägheit weiter, und um sich zu bewegen, braucht sie keine Kraft. Auf die Kiste wirken nur die Gewichtskraft, die Normalkraft und die Gleitreibung. Die Schritte 2 und 3 stimmen.

\abc
Entlang des Bodens wirkt nur die Gleitreibung, gegen die Bewegung (in Bewegungsrichtung positiv):
\[
 -F_\mathrm{R} = m\,a
\]
Die Beschleunigung zeigt also gegen die Bewegung, die Kiste wird langsamer. Ihr Betrag:
\begin{align}
 a &= \aKF = \frac{\FR}{\mK} = \aK \approx \result{\aKP{0}{2}}
\end{align}
gegen die Bewegungsrichtung. (Weil $F_\mathrm{R} = \mu\,m\,g$ ist, gilt $a = \mu\,g$: Die Masse spielt keine Rolle.)

\abc
Die Kiste bremst gleichmässig von $v_0 = \vNO$ auf null ab: $v_0^2 = 2\,a\,s$.
\begin{align}
 s &= \sKF \\
   &= \frac{\left(\vN\right)^2}{2\cdot\aK} \\
   &= \sK \approx \result{\sKP{0}{2}}
\end{align}
Die Kiste bleibt stehen, weil die Reibung sie abbremst, bis die Geschwindigkeit null ist, und nicht, weil eine Bewegungskraft „aufgebraucht“ ist.
\end{abcliste}
